\section{Three forms with matching coefficient columns}\label{sec:forms}

\subsection{Designing one function for all three grids}

To obtain the coefficient pattern in \eqref{eq:preview}, we evaluate
a single rational function $R_n$ on each grid $j/b$, for $b=1,2,3$,
and sum with the alternating signs and odd indices of the beta series.

Write these signs as weights $\ch(j)$:
$\ch(2k)=0$ and $\ch(2k+1)=(-1)^k$ for integers $k\geq0$.
Thus the weights are $1,0,-1,0,1,0,-1,0,\ldots$ for $j=1,2,\ldots$,
and
\[
 \sum_{j=1}^{\infty}\ch(j)R_n(j/b)
   =R_n(1/b)-R_n(3/b)+R_n(5/b)-\cdots.
\]
The weights have period $4$. We therefore space the poles by $4$:
a partial-fraction term at $-4h$ satisfies
\[
 \frac{1}{(j/b+4h)^i}=\frac{b^i}{(j+4bh)^i}.
\]
Since $4bh$ is a multiple of $4$, the shift preserves the weights.
Poles at $0,-4,\ldots,-4n$ therefore give the required factor $b^i$.

The numerator serves two purposes. Reflection about $-2n$ pairs the
poles; we arrange that $R_n(-4n-t)=R_n(t)$ for even $n$, so the paired
coefficients cancel at every odd order. We also make the initial terms
of each alternating sum vanish, leaving only a tail. If the last
positive zero is $rn$, its reflection is $-(r+4)n$. The blocks of zeros
must therefore span an interval of length $(2r+4)n$.

The three grids have fractional parts $0,1/2,1/3,2/3$.
We use the representatives
\[
 \Theta=\{1,\tfrac12,\tfrac13,\tfrac23\},
\]
where $1$ represents the integer class. We choose $r=7$, giving a block
length of $18n$, and take the largest beta index to be $364$.
Longer zero blocks remove more initial terms but increase the numerator
degree. Higher pole orders improve decay but enlarge both the list of
beta values and the multiplier needed to clear denominators.
Sections~\ref{sec:decay} and~\ref{sec:certificate} verify that our
choices give sufficient decay and the required nonvanishing.

For positive even $n$, the numerator with these zero blocks is
\begin{equation}
 P_n(t)=(t-7n)\prod_{\theta\in\Theta}\prod_{j=0}^{18n-1}
                         (t-7n+j+\theta).\label{eq:pn}
\end{equation}
The factor $t-7n$ adds the right endpoint to the integer block.
The integer roots of $P_n$ are exactly $-11n,-11n+1,\ldots,7n$;
the half-integer block reflects into itself, and the third-shift blocks
exchange. In particular, $P_n$ has a simple zero at each proposed
pole $-4h$. To leave pole order $364$ after cancellation, we use
the denominator $\prod_{h=0}^n(t+4h)^{365}$.

\subsection{Choosing the arithmetic normalization}

Multiplication by a constant preserves the zeros, poles and reflection
symmetry. We choose this constant to control the denominators of the
partial-fraction coefficients. For each numerator block, we divide by
$(n!)^{18}$, using the divisibility of products of $n$ consecutive terms.
There are four blocks, hence the denominator $(n!)^{72}$.
The fractional shifts require compensating powers of $2$ and $3$,
supplied by $K^n$ with $K=2^{36}3^{54}$.
Section~\ref{sec:arithmetic} proves the resulting divisibility bounds.

Expansion at a pole $-4h$ introduces further denominators.
The constant terms of the other pole factors have product
\[
 \prod_{\substack{0\leq k\leq n\\k\ne h}}(4(k-h))^{365}
 =4^{365n}(-1)^{365h}\bigl(h!(n-h)!\bigr)^{365}.
\]
Multiplication by $4^{365n}$ cancels the power of $4$.
Multiplication by $((n+1)!)^{365}$ then leaves the integer factor
$(-1)^{365h}((n+1)\binom nh)^{365}$. We therefore set
\begin{equation}
 C_n=4^{365n}K^n\frac{((n+1)!)^{365}}{(n!)^{72}},\qquad
 R_n(t)=\frac{C_nP_n(t)}{\prod_{h=0}^{n}(t+4h)^{365}}.
 \label{eq:rn}
\end{equation}
The size of $C_n$ must be included in the decay estimate of
Section~\ref{sec:decay}. The numerator has degree $72n+1$, so $R_n$ has degree at
infinity $-293n-364$. Each pole has order exactly $364$. Write
\begin{equation}\label{eq:partials}
 R_n(t)=\sum_{h=0}^{n}\sum_{i=1}^{364}
              \frac{a_{i,h}(n)}{(t+4h)^i},\qquad
 \rho_i(n)=\sum_{h=0}^{n}a_{i,h}(n).
\end{equation}

\begin{samepage}
\begin{lemma}\label{lem:forms}
For $b=1,2,3$, define
\begin{align}
 r_{n,b}&=\sum_{j=1}^{\infty}\ch(j)R_n(j/b),\label{eq:rnb}\\
 q_b(n)&=-\sum_{h=0}^{n}\sum_{i=1}^{364}a_{i,h}(n)b^i
                   \sum_{j=1}^{4bh}\frac{\ch(j)}{j^i}.
                   \label{eq:qb}
\end{align}
Then
\begin{equation}\label{eq:forms}
 r_{n,b}=q_b(n)+\sum_{\substack{2\leq i\leq364\\i\text{ even}}}
                         b^i\rho_i(n)\beta(i).
\end{equation}
Moreover,
\begin{equation}\label{eq:positive-bound}
 |r_{n,b}|\leq\sum_{j>7bn}R_n(j/b).
\end{equation}
\end{lemma}
\end{samepage}

\begin{proof}
The roots of $P_n$ are symmetric about $-2n$. The integer and half-integer
blocks each reflect into themselves, and the one-third and two-thirds
blocks exchange. Because $n$ is even, numerator and denominator change
by the same sign under $t\mapsto-4n-t$. Therefore
\[
 R_n(-4n-t)=R_n(t),\qquad
 a_{i,n-h}(n)=(-1)^i a_{i,h}(n),\qquad \rho_i(n)=0\quad(i\text{ odd}).
\]
For each partial fraction, shifting the summation variable by $4bh$
preserves $\ch$ and gives
\[
 \sum_{j=1}^{\infty}\frac{\ch(j)}{(j/b+4h)^i}
 =b^i\left(\beta(i)-\sum_{j=1}^{4bh}\frac{\ch(j)}{j^i}\right).
\]
Here $\beta(1)$ denotes its convergent alternating series. The $i=1$
sums converge because the partial sums of the weights are bounded;
there are only finitely many partial-fraction sums. The sum of the
$\beta(1)$ coefficients is zero, as is the sum at every other odd index.
This proves \eqref{eq:forms}.
The original series \eqref{eq:rnb} is absolutely convergent.

Every positive grid point $j/b\leq7n$ is a numerator zero. Beyond $7n$,
all factors of $R_n$ are positive. Taking absolute values now gives
\eqref{eq:positive-bound}. We use this positive tail only to bound size;
nonvanishing of the alternating sum requires a separate argument.
\end{proof}
